\section{Discussion}
\label{sec:discussion}

\begin{table*}[t]
\centering
\caption{Trust and cost comparison. Bytes: encoder URL token, 60-character path. Verify: native OpenSSL~3.5.4 medians. Dissemination: simulated fleet traffic for $E=500$ ($\bar\lambda=2000$/s) or per-session staple bytes.}
\label{tab:trust}
\scriptsize
\setlength{\tabcolsep}{3pt}
\begin{tabular}{@{}llp{2.9cm}p{2.0cm}rp{3.0cm}p{2.6cm}@{}}
\toprule
Scheme & Secret at edge & One compromised edge mints for & PQ assumption & Bytes & Verify (\textmu s) & Dissemination \\
\midrule
B1 HMAC & yes (key lifetime) & all edges, all paths of the key & unconditional & 166 & 0.30 & none \\
B10 multi-MAC, $k{=}3$ & yes & its own domain & unconditional & 187 & 0.30 & none \\
B12 KEM-wrapped HMAC & yes (epoch keys) & all edges, live epochs & unconditional & 171 & 0.30 & 283.5~MB/s key push \\
B2 Ed25519 & no & none & none (Shor) & 209 & 94.6 & none \\
B11 delegated Ed25519 & no & none & no break within $L+\vmax$ & 214 & 94.6 per (key, edge) & 0.383~MB/s ($L=1$~h) \\
B5 ML-DSA-44 & no & none & unconditional & \num{3350} & 100.9 per request & none \\
B5v + verify cache & no & none & unconditional & \num{3350} & 100.9 first, 2.19 hit & none \\
\textbf{M2} stapled ML-DSA-44 & no & none & unconditional & 148 & 100.9 per (token, edge), $\le$0.21 hit & staple \num{2921}~B/session (h2) \\
EMBER pooled, staple-all & no & none & unconditional & 268 & 0.89 path ($d{=}5$) + 100.9 per (root, edge) & staple \num{2968}~B/session (h2) \\
EMBER per-tenant, pull & no & none & unconditional & 161--247 & 0.24--0.75 path ($d{=}0$--4) + root & 8.41~MB/s ($\Delta=1$~s) \\
\bottomrule
\end{tabular}
\end{table*}

\subsection{Deployment Guidance}
Table~\ref{tab:trust} summarizes the trust and cost trade-offs. Our recommendations follow from the measurements.

\textbf{1. If edges may hold a minting key, use HMAC.} HMAC (B1) is post-quantum and the cheapest option. Per-domain multi-MAC (B10) confines a multi-CDN breach to one vendor at HMAC cost.

\textbf{2. Otherwise, staple.} Detaching the signature from the token and sending it once per edge connection fixes the edge for any signature scheme. M2 matches HMAC in edge CPU (Table~\ref{tab:p1}) and is the smallest post-quantum option in URL form; in cookie form it is within 0.4\% of the smallest (Table~\ref{tab:wire}). Clients should staple on every new edge contact, not only on the first request (Fig.~\ref{fig:pmove}). Stapling also removes the root service as an availability dependency.

\textbf{3. Batch only when the issuer is the bottleneck, and then pool.} Per-token signing costs the issuer $\lambda\times408.8$\us, about 0.41 cores per $10^3$ tokens/s per key. Pooled EMBER at $\Delta=10$~ms costs at most $1/\Delta$ root signatures per second plus hashing. In our runs this stayed at or below 0.057 cores up to \num{20000} tokens/s, and at $\lambda\ge10^3$/s trees hold at least about 10 tokens. Below about $10^3$ tokens/s per issuer key, the saving is under half a core: 0.046 cores at 200/s. That does not justify 5.4~ms of extra delay, 120--184 more bytes per token and the extra path hashing at the edge, so M2 is the better choice. Above that rate, use one pooled tree per issuer key with $\Delta\approx10$~ms and staple-all. Pooling requires that the tenants of a pool trust a shared issuer key. With per-tenant keys, batching saves only 1.2$\times$, and M2 should be used. Work-conserving closing self-limits issuer CPU to one core per pool at sub-millisecond delay (0.62~ms at \num{20000}/s), which suits pools whose rate fits on one signer core.

\textbf{4. Consider B11 if its assumption holds.} Delegated Ed25519 is smaller and cheaper to disseminate. It is appropriate when breaking one Ed25519 key is believed to take longer than $L+\vmax$ (Section~\ref{sec:res-b11}). Short $\vmax$ is the effective lever.

\textbf{5. Operational settings.} Use a node size of $n=16$ for category-1 tokens and $n=32$ above that. Add a per-$\mathit{kid}$ public seed to the tweak. Set $\vmax=1$~h and a skew tolerance $\tau=5$~s. Apply per-client limits on staples and fetches, and fetch roots only for configured $\mathit{kid}$s and recent epochs. Do not deploy per-tenant push or rely on roll-up to reduce misses.

\subsection{Limitations}
\label{sec:limits}
\begin{itemize}
\item \emph{One VM.} All prototype components share one 4-vCPU KVM guest. Network namespaces and a userspace link emulator stand in for \texttt{netem}, and there is no physical NIC and no TLS. P1 runs over loopback, so it measures edge CPU per request, not network effects. Separate P1 campaigns differed by about 10\% in absolute value.
\item \emph{Two native runs.} Absolute costs differ by 15--20\% across a VM reboot. Everything except the controller study and the simulator consistency check uses the second run (Appendix~\ref{app:secondary}).
\item \emph{Simulator scope.} The simulator is validated against the prototype only for miss and byte counts, at small scale ($V\le50$), and without re-routes or a shield tier. It has no edge-CPU or link queueing, so it under-predicts pull-miss latency for timer-aligned bursts (118--137~ms measured vs.\ 75.9~ms). Workload parameters (Zipf exponents, RTTs, $K$, $p_{\text{move}}$) are assumptions, and $p_{\text{move}}$ alone determines staple-first waiting.
\item \emph{Clients.} We tested Chrome~153, curl/nghttp2 and Go, but not Firefox or Safari.
\item \emph{Security.} We give a proof sketch, not a machine-checked proof. The public-seed variant was not implemented. Bearer-token replay is unchanged from the status quo. Pooling concentrates minting in one key.
\end{itemize}
